\section{Introduction and main results}

For $N\ge8$, the cycle square $C_N(1,2)$ has vertex set $\Z/N\Z$ and
edges $\{i,i+1\}$ and $\{i,i+2\}$. A real signing assigns $\pm1$ to each
edge. For its symmetric adjacency matrix $A$, write
\[
 \rho(A)=\max\{|\lambda|:\lambda\in\spec(A)\},
 \qquad \mu_N=\min_\sigma\rho(A_\sigma).
\]
The minimization controls both spectral edges. The benchmark in the
original twisted-optimality question \cite{Suvagiya2026old} is
\begin{equation}
 \tw(N)^2=4+2\cos\frac{2\pi}{N}+2\cos\frac{4\pi}{N}
 \qquad(N\text{ even},\ N\ge8).
 \label{eq:twisted}
\end{equation}

Our main result is a uniform bound for a family whose local structure is
fixed but whose cell lengths and cell count are arbitrary. Set
\begin{equation}
 t=(1,1,-1,1,-1,-1,1,-1),\qquad
 w_j=t^j\mathbin{\|}(1,-1),\qquad h_j=8j+2\quad(j\ge1).
 \label{eq:word}
\end{equation}
Here $\|$ denotes concatenation. For a length-$N$ triangle-sign word
$\tau$ and $\alpha\in\{\pm1\}$, define $A_N(\tau,\alpha)$ by the edge signs
\begin{equation}
 a_i=\begin{cases}1&i<N-1,\\\alpha&i=N-1,\end{cases}
 \quad
 \sigma\{i,i+1\}=a_i,\qquad
 \sigma\{i,i+2\}=\tau_i a_i a_{i+1},
 \label{eq:gauge-definition}
\end{equation}
with cyclic indices. The triangle signs are $\tau_i$, and the
Hamilton-cycle sign product is $\alpha$.
Define the two rational caps
\[
 \czero=\frac{198}{25}=7.92,\qquad
 \cone=\frac{790537}{100000}=7.90537.
\]

\begin{theorem}[Unequal-cell bound]\label{thm:cells}
Let $r\ge1$, $j_0,\ldots,j_{r-1}\ge1$, and
$\tau=w_{j_0}\|\cdots\|w_{j_{r-1}}$, of total length $N$.
For either holonomy $\alpha$,
\begin{align}
 \min_i h_{j_i}\ge106&\quad\Longrightarrow\quad
       \rho(A_N(\tau,\alpha))^2<\czero,\label{eq:cell-cap0}\\
 \min_i h_{j_i}\ge202&\quad\Longrightarrow\quad
       \rho(A_N(\tau,\alpha))^2<\cone.\label{eq:cell-cap1}
\end{align}
The constants do not depend on $r$, and the lengths need not be equal.
\end{theorem}

The allowed cell words are part of the hypotheses. The theorem does not
cover arbitrary signings or arbitrary arrangements of local defects.
Its proof retains six coordinates at each cell boundary and eliminates the
four-site interior chains exactly. A common positive limiting core and a
quadratic-form estimate with two chain incidences per boundary control the
resulting $6r$-dimensional cyclic matrix. The error is independent of $r$;
no count of localized eigenmodes is assumed.

The single-cell family permits a sharper finite completion.
\begin{theorem}[One-cell bound]\label{thm:one-cell}
For every $k\ge1$,
\[
 \rho(A_{8k+2}(w_k,+1))^2<\cone.
\]
Moreover,
\begin{equation}
 \frac{7905369}{1000000}
 <\sup_{k\ge1}\rho(A_{8k+2}(w_k,+1))^2
 \le\frac{790537}{100000}.
 \label{eq:supremum}
\end{equation}
\end{theorem}
The interval in \eqref{eq:supremum} has width $10^{-6}$. Its upper endpoint
is non-strict; no convergence, monotonicity, or exact supremum is asserted.

To cover each nonzero even residue modulo eight, put
\[
 W_{k,r}=\mathop{\|}_{a=0}^{r-1}w_{\lfloor(k+a)/r\rfloor},
 \qquad r\in\{1,2,3\},\ k\ge6.
\]
The cell parameters sum to $k$, so the word has length $8k+2r$.
\begin{theorem}[Explicit finite competitors]\label{thm:witnesses}
For $r\in\{1,2,3\}$ and $k\ge6$, the matrix
\[
 B_{k,r}=A_{8k+2r}(W_{k,r},(-1)^{r+1})
\]
satisfies
\begin{equation}
 \rho(B_{k,r})^2<\czero<\tw(8k+2r)^2.
 \label{eq:nonzero-witnesses}
\end{equation}
For $r=1$ the stronger cap $\cone$ holds. Together with the period-eight
family, these constructions prove
\[
 \mu_N<\tw(N)
 \quad\text{for }N\in\{32,40\}\cup\{N\ge48:N\text{ even}\}.
\]
\end{theorem}

This witness range already occurs in the earlier Target A project
\cite{EarlierTargetA}. The contribution here is a common transparent
analytic mechanism, sharper quantitative bounds, and a finite completion
with directly specified signings. For odd $k$ in residue four, and for
$3\nmid k$ in residue six, our placements need not be the older balanced-gap
placements. Their finite certificates are recomputed for the words above.
We do not reprove the universal lower bounds at complementary smaller
orders, determine $\mu_N$ when the twisted signing fails, or classify
minimizing signings.

The literature target has also changed. The original preprint
\cite{Suvagiya2026old} was withdrawn and merged into
\cite{Suvagiya2026}. Its revised Theorem 26 supplies the positive-holonomy
period-eight value
\begin{equation}
 \etaedge=4+\sqrt{10+2\sqrt5},\qquad \rstar=\sqrt{\etaedge},
 \label{eq:rstar}
\end{equation}
for $N=8m$, $m\ge4$, while Conjecture 28 proposes $\mu_{8m}=\rstar$
throughout that range. The exact negative-holonomy formula from the earlier
project manuscript \cite{EarlierPeriodEight} is strictly smaller at every
finite size. We provide its full proof and apply it to that revised claim.
This application is distinct from a priority claim for the inherited formula.

The use of holonomy belongs to the established character approach to
spectral graph theory; Luo and Roy \cite[Theorems 3.6 and 3.13]{LuoRoy2026}
study character families and endpoint attainment. Those results do not
supply the fixed-support minimum here. Likewise, signed moments averaged
over all signings, as in Chen, van Dam, and Bu
\cite[Theorem 3.1]{ChenDamBu2024}, are not pointwise lower bounds for every
signing. Our estimates instead arise from an explicit fiber calculation
and exact positive-definiteness tests after a uniform analytic reduction.
The latter proofs are finite-certificate-assisted, with independent exact
replays; their finite checks are specified in Section~\ref{sec:finite}.
